\sectiontheme{rustaccent}
\section{Exports: Markdown, OpenAPI, Postman}
\label{sec:exports}

Three commands don't emit source code at all --- they \emph{describe} a
module in a standard interoperable format, reading the same definition
that produces the Go/Swift/Kotlin/JS clients, so none of them can drift
away from the real implementation.

\subsection{\texttt{emi md} --- Markdown docs}
The fastest way to get human-readable, always up to date documentation
for a module: a title, description, table of contents, and one section
per action listing method, URL, CLI name, and request/response fields.
\begin{lstlisting}
emi md --path module.yaml --output ./out
# ./out/postsModule.md
\end{lstlisting}
Output for a single-action module looks like:
\begin{lstlisting}
### `GetSinglePostAction`
| | |
|---|---|
| **Method** | `GET` |
| **URL** | `/posts/:id` |
| **CLI** | `get-single-post` |
#### Response
| Field | Type |
|---|---|
| `userId` | `int64` |
\end{lstlisting}
Committed next to the module, served in a docs site, or rendered straight
into a README.

\subsection{\texttt{emi openapi} --- OpenAPI 3}
A \emph{describer}, not a compiler: reads the definition, produces a
standard \href{https://spec.openapis.org/oas/v3.0.3}{OpenAPI 3} document
for Swagger UI, Redoc, code generators, or API gateways.
\begin{lstlisting}
emi openapi --path module.yaml --output ./out --json
# ./out/postsModule.openapi.yaml (+ .json)
\end{lstlisting}
Route parameter syntax is rewritten (\texttt{/posts/:id} becomes
\texttt{/posts/\{id\}}), path parameters are promoted, and Emi field
types map onto the correct OpenAPI \texttt{type}/\texttt{format} pairs
(e.g.\ \texttt{int64} $\to$ \texttt{type: integer, format: int64}).

\subsection{\texttt{emi postman} --- Postman v2.1 collection}
Every action becomes a Postman request with method, URL, and the standard
Emi headers pre-filled, ready to import and start hitting endpoints
immediately.
\begin{lstlisting}
emi postman --path module.yaml \
  --host localhost --port 8080 --output ./out
\end{lstlisting}
Host and port are emitted as Postman environment variables
(\texttt{\{\{HOST\}\}}, \texttt{\{\{PORT\}\}}) so switching between
local, staging and production is a matter of changing the active
environment, and the common Emi headers (\texttt{authorization},
\texttt{role-id}, \texttt{workspace-id}) are wired to
\texttt{\{\{AUTH\}\}}/\texttt{\{\{RID\}\}}/\texttt{\{\{WID\}\}}
variables, filled once and reused across every request.

\begin{pullquote}
All three exporters share the same two flags: \texttt{--path} (the
definition file) and \texttt{--output} (where to write; omit it and the
result prints to stdout so it can be piped or inspected). Unlike the
code-generating compilers (\S\ref{sec:additional-targets}), none of the
three currently registers its own \texttt{--tags} --- there is nothing
yet to toggle in a document that's just read back off the definition.
\end{pullquote}
